\documentclass[11pt]{article}
\usepackage[margin=1in]{geometry}
\usepackage{amsmath,amssymb,amsthm}
\usepackage{graphicx}
\usepackage{hyperref}

\newtheorem{theorem}{Theorem}
\newtheorem{corollary}{Corollary}
\newtheorem{proposition}{Proposition}
\newcommand{\N}{\mathbb N}

\title{Deficit classifications for diagonals of square partial isometries}
\author{Major partial answer to Question 6.2 of arXiv:1602.08435}
\date{12 August 2026}

\begin{document}
\maketitle

\noindent\textbf{Status.} \texttt{candidate\_major\_partial\_likely\_valid}.

\begin{abstract}
Jasper, Loreaux, and Weiss ask for the diagonals of partial isometries with
prescribed kernel and cokernel dimensions, in particular square partial
isometries.  For a sequence $d=(d_j)$, put
\[
 \Delta(d)=\sum_j(1-|d_j|)\in[0,\infty].
\]
We prove that every diagonal of a square partial isometry of defect $r$
satisfies $\Delta(d)\ge r$, with the extended interpretation when
$r=\infty$.  This bound is sharp.  It gives a complete classification for
infinite defect: precisely $|d_j|\le1$ and $\Delta(d)=\infty$.  In fact,
under infinite deficit every prescribed square defect
$r=0,1,2,\ldots,\infty$ is realizable.  For finite $r$, we also prove the
converse when $\Delta(d)=r$; for positive finite $r$ the converse holds as
well whenever only finitely many entries fail to be unimodular.  All
constructions preserve real scalars for real diagonals.
\end{abstract}

\section{Source question and notation}

Let $H$ be a separable infinite-dimensional Hilbert space with fixed
orthonormal basis $(e_j)_{j\in\N}$.  A partial isometry $V$ is
\emph{square of defect $r$} when
\[
 \dim\ker V=\dim\ker V^*=r,
 \qquad r\in\{0,1,2,\ldots,\infty\}.
\]
Its diagonal is $d_j=\langle Ve_j,e_j\rangle$.  Question 6.2 of the source
asks for these diagonals; see Figure~\ref{fig:source}.

We use two results already stated in the source paper \cite{source}:
finite-dimensional Thompson's theorem and its unitary-diagonal theorem.  The
former says that $z_1,\ldots,z_N$ are the diagonal of a matrix with singular
values $s_1\ge\cdots\ge s_N$ exactly when the Thompson inequalities hold.
The latter says that $z=(z_j)$ is the diagonal of a unitary exactly when
\[
 |z_j|\le1,
 \qquad
 2\bigl(1-\inf_j|z_j|\bigr)
 \le \sum_j(1-|z_j|).
\tag{1}\label{eq:unitary}
\]

\section{Idea of the proof}

There are two complementary mechanisms.  For necessity, look simultaneously
at the kernel and cokernel projections.  Their diagonal masses cannot both be
large at a coordinate where the diagonal entry of the partial isometry is
nearly unimodular.  Cauchy--Schwarz makes this quantitative and summing recovers
the defect.

For sufficiency, total modulus deficit is a budget for zero singular values.
On a finite block, one unit of deficit buys one zero singular value through
Thompson's theorem.  Infinite deficit can therefore be divided into finite
blocks to create infinite defect, or one finite block can create any prescribed
finite defect while a unitary realizes the remaining infinite-deficit tail.
At exact finite equality, the entire deficit is instead realized as the
diagonal of a finite-rank projection, whose complement becomes the initial
projection of the desired partial isometry.

\section{The sharp necessary inequality}

\begin{proposition}[Two-sided defect bound]\label{prop:necessary}
Let $V$ be a partial isometry with diagonal $d$.  Put
\[
 P=I-V^*V,\quad Q=I-VV^*,\quad
 p_j=\langle Pe_j,e_j\rangle,\quad
 q_j=\langle Qe_j,e_j\rangle.
\]
Then, pointwise,
\[
 p_j+q_j\le 2(1-|d_j|).                                      \tag{2}\label{eq:pointwise}
\]
Consequently, if $V$ is square of finite defect $r$, then
\[
 |d_j|\le1\quad\hbox{and}\quad \Delta(d)\ge r.              \tag{3}\label{eq:necessary}
\]
If both defects are infinite, then $\Delta(d)=\infty$.
\end{proposition}

\begin{proof}
The vectors $Ve_j$ and $VV^*e_j$ lie in the final space of $V$, and
\[
 d_j=\langle Ve_j,e_j\rangle
     =\langle Ve_j,VV^*e_j\rangle.
\]
Moreover,
\[
 \|Ve_j\|^2=1-p_j,
 \qquad
 \|VV^*e_j\|^2=\|V^*e_j\|^2=1-q_j.
\]
Cauchy--Schwarz followed by the arithmetic--geometric mean inequality gives
\[
 |d_j|
 \le\sqrt{(1-p_j)(1-q_j)}
 \le 1-\frac{p_j+q_j}{2},
\]
which is \eqref{eq:pointwise}.  If $P,Q$ have rank $r<\infty$, Tonelli's
theorem and the trace formula for positive finite-rank operators yield
\[
 2r=\sum_j(p_j+q_j)
 \le 2\sum_j(1-|d_j|).
\]
If both ranks are infinite, both projection-diagonal sums diverge, and the
same inequality forces $\Delta(d)=\infty$.  Contractivity of $V$ gives
$|d_j|\le1$.
\end{proof}

The use of both defect projections is essential: either one alone gives only
the weaker lower bound $r/2$.

\section{Finite Thompson blocks}

\begin{proposition}[Defect block]\label{prop:block}
Let $B$ be a finite index set of size $N$, let $1\le r\le N$, and
suppose
\[
 |d_j|\le1\quad(j\in B),
 \qquad \sum_{j\in B}(1-|d_j|)\ge r.                         \tag{4}\label{eq:block}
\]
Then there is an $N$-dimensional partial isometry of defect $r$ whose
diagonal is $(d_j)_{j\in B}$.  It can be chosen real when the $d_j$ are
real.
\end{proposition}

\begin{proof}
Apply finite Thompson to the singular-value list
\[
 s=(\underbrace{1,\ldots,1}_{N-r},
       \underbrace{0,\ldots,0}_{r}).
\]
Write $b_1\ge\cdots\ge b_N$ for the decreasing rearrangement of the
$|d_j|$.  For $k\le N-r$, $\sum_{i\le k}b_i\le k=\sum_{i\le k}s_i$.
For $k>N-r$,
\[
 \sum_{i\le k}b_i\le\sum_{i\le N}b_i
 =N-\sum_{j\in B}(1-|d_j|)\le N-r=\sum_{i\le k}s_i.
\]
The final Thompson inequality follows from
\[
 \sum_{i=1}^{N-1}b_i-b_N\le\sum_{i=1}^N b_i\le N-r
=\sum_{i=1}^{N-1}s_i-s_N.
\]
Thus Thompson gives a matrix
$A$ with diagonal $(d_j)$ and singular values $s$.  Since the spectrum
of $A^*A$ is contained in $\{0,1\}$, $A^*A$ is a projection, so $A$
is a partial isometry.  Both defects equal $r$.  Thompson's real conclusion
gives the last assertion.
\end{proof}

\section{Infinite deficit: a complete classification}

\begin{theorem}[Defect flexibility]\label{thm:flex}
Suppose $|d_j|\le1$ and $\Delta(d)=\infty$.  Then, for every
\[
 r\in\{0,1,2,\ldots,\infty\},
\]
there is a square partial isometry of defect $r$ with diagonal $d$.
If $d$ is real, it can be realized on a real Hilbert space.
\end{theorem}

\begin{proof}
For $r=0$, condition \eqref{eq:unitary} is automatic, so the source's
unitary theorem applies.

Fix $1\le r<\infty$.  Choose a finite set $B$ with
$\sum_{j\in B}(1-|d_j|)\ge r$.  Proposition~\ref{prop:block} gives a
defect-$r$ partial isometry on the corresponding finite-dimensional block.
Removing finitely many terms leaves infinite deficit, so the complementary
sequence satisfies \eqref{eq:unitary} and is the diagonal of a unitary.  The
orthogonal direct sum has both defects $r$.

For $r=\infty$, greedily partition $\N$ into finite sets
$B_1,B_2,\ldots$ with
\[
 \sum_{j\in B_k}(1-|d_j|)\ge1.
\]
This is possible because every tail after finitely many removals still has
infinite sum.  Apply Proposition~\ref{prop:block} with $r=1$ on every
block.  Their direct sum has infinite-dimensional kernel and cokernel.  All
ingredients have real versions.
\end{proof}

Combining necessity and sufficiency gives the promised full subcase.

\begin{corollary}[Infinite square defect]\label{cor:infinite}
A complex sequence $d$ is the diagonal of a square partial isometry with
infinite-dimensional kernel and cokernel if and only if
\[
 |d_j|\le1\quad\hbox{for all }j,
 \qquad \sum_j(1-|d_j|)=\infty.
\]
The same characterization holds over the reals for real $d$.
\end{corollary}

\section{Two exact finite-defect regimes}

\begin{theorem}[Equality boundary]\label{thm:equality}
Let $r<\infty$.  If $|d_j|\le1$ and $\Delta(d)=r$, then $d$ is the
diagonal of a square partial isometry of defect $r$.
\end{theorem}

\begin{proof}
Set $a_j=1-|d_j|$.  The sequence $a=(a_j)$ lies in $[0,1]$ and has sum
$r$.  By the finite-rank Schur--Horn theorem (equivalently, the finite-trace
case of Kadison's carpenter theorem stated in the source), there is a
rank-$r$ projection $P$ with diagonal $a$.  Thus $I-P$ has diagonal
$(|d_j|)$.  Choose a diagonal unitary $W=\operatorname{diag}(w_j)$ with
$w_j|d_j|=d_j$, and put
\[
 V=W(I-P).
\]
Then $V^*V=I-P$ and $VV^*=W(I-P)W^*$; both complementary projections have
rank $r$, and the diagonal of $V$ is $d$.  For real $d$, the source's
real finite-rank Schur--Horn theorem gives a real $P$, and
$w_j\in\{-1,1\}$ may be used.
\end{proof}

\begin{theorem}[Eventually unimodular sequences]\label{thm:eventual}
Fix $1\le r<\infty$, and suppose $|d_j|\le1$ with $|d_j|=1$ except at
finitely many indices.  Then $d$ is the diagonal of a square partial
isometry of defect $r$ if and only if $\Delta(d)\ge r$.
\end{theorem}

\begin{proof}
Necessity is Proposition~\ref{prop:necessary}.  For sufficiency, let $B$
contain all indices with $|d_j|<1$.  Proposition~\ref{prop:block} realizes
that finite block with defect $r$; the complementary diagonal is itself a
diagonal unitary.  Their direct sum proves the claim.
\end{proof}

\section{Scope, novelty, and remaining obstruction}

The complete open question is not claimed.  The unresolved regime after these
results is finite $r$ with
\[
 r<\Delta(d)<\infty
\]
and infinitely many non-unimodular entries.  Finite-section partial
isometries need not have a strongly convergent subsequence: off-diagonal
unitary mass can escape even when the finite-rank kernel and cokernel
projections are tight.  Nor can the tail always be delegated to a unitary;
superlacunary deficit sequences can violate \eqref{eq:unitary} after every
finite removal.  These are genuine gaps, not merely omitted details.

Cheap-index searches and bounded arXiv searches through 12 August 2026 found
no later solution of Question 6.2 and no statement of
Corollary~\ref{cor:infinite}.  Novelty confidence is moderate: the arguments
use classical ingredients and could be unindexed folklore.  Human review
should focus on Proposition~\ref{prop:block}, the projection-diagonal use in
Theorem~\ref{thm:equality}, and the claim that the selected question remains
open in the literature.

\begin{thebibliography}{9}
\bibitem{source}
J. Jasper, J. Loreaux, and G. Weiss,
\emph{Thompson's theorem for compact operators and diagonals of unitary
operators}, arXiv:1602.08435 (2016).
\end{thebibliography}

\clearpage
\begin{figure}[p]
\centering
\includegraphics[width=\textwidth]{figures/open_question_crop.png}
\caption{Question 6.2 on source page 21 asks for the diagonals of partial
isometries with prescribed kernel and cokernel dimensions, especially square
partial isometries.}
\label{fig:source}
\end{figure}

\end{document}
